\documentclass[UTF8]{ctexart}
\usepackage{xeCJK}
\usepackage{geometry}
\geometry{left=2cm,right=2cm,top=2.5cm,bottom=2.5cm}
\setlength{\headheight}{12.7pt}

\usepackage{titlesec}
\titleformat{\section}
  {\normalfont\large\bfseries}
  {\thesection}
  {1em}
  {}
\titlespacing*{\section}
  {0pt}
  {3.5ex plus 1ex minus .2ex}
  {2.3ex plus .2ex}

\usepackage{amsmath}
\usepackage{amssymb}
\usepackage{amsthm}
\usepackage{enumitem}
\setlist[enumerate]{itemsep=0pt,topsep=0pt,parsep=0pt,leftmargin=0pt,itemindent=2.5em}
\setlist[itemize]{itemsep=0pt,topsep=0pt,parsep=0pt,leftmargin=0pt,itemindent=2.5em}
\usepackage{fancyhdr}
\pagestyle{fancy}
\fancyhf{}
\usepackage{graphicx}
\usepackage{hyperref}
\usepackage{indentfirst}
\usepackage{lmodern}

\begin{document}
\setlength{\abovedisplayskip}{1pt}
\setlength{\belowdisplayskip}{1pt}

\section{生成子结构}

\hypertarget{sec:定义1.1}{}
\noindent\textbf{定义1.1 (生成子结构)}

给定一阶语言$\mathcal{L}$, $\mathcal{L}$-结构$B$和非空的$\Omega\subset B$.
在$B$中定义:
\begin{enumerate}
    \item 初始集$\Omega\cup\{c^B\mid c\in\mathcal{C}\}$,
    \item 归纳规则: 对每个函数符号$f\in\mathcal{F}$, 指定映射$f^B$.
\end{enumerate}
从$\Omega$和归纳规则可以归纳定义$\langle\Omega\rangle\subset B$,
称为$\Omega$在$B$中生成的子结构.

\vspace{10pt}

\hypertarget{sec:定理1.1}{}
\noindent\textbf{定理1.1}

给定一阶语言$\mathcal{L}$, $\mathcal{L}$-结构$B$和非空的$\Omega\subset B$,
则$\langle\Omega\rangle$是$B$的包含$\Omega$的
最小的\href{03 子结构#nameddest=sec:定义1.1}{子结构}.

\noindent{\kaishu 证明.}

首先依定义可知$\langle\Omega\rangle$含有所有的$c^B$并且对所有$f^B$封闭, 所以它是子结构.

其次, 假设$A\supset\Omega$也是$B$的子结构, 那么:
\begin{enumerate}
    \item 对任意的$b\in\Omega\cup\{c^B\mid c\in\mathcal{C}\}$, 都有$b\in A$,
    \item 对任意的$n$元函数符号$f$, 假设$a_1,\cdots,a_n\in A$, 那么
    $f^B(a_1,\cdots,a_n)\in A$.
\end{enumerate}

所以归纳证明得: 对任意的$b\in\langle\Omega\rangle$都有$b\in A$,
从而$\langle\Omega\rangle\subset A$. \hfill $\qedsymbol$

\vspace{10pt}

\hypertarget{sec:定理1.2}{}
\noindent\textbf{定理1.2}

给定一阶语言$\mathcal{L}$, $\mathcal{L}$-结构$B$和非空的$\Omega\subset B$,
$|\langle\Omega\rangle|\leqslant\max\{|\Omega|,|\mathcal{L}|\}$.

\noindent{\kaishu 证明.}

显然$|\Omega\cup\{c^B\mid c\in\mathcal{C}\}|, |\mathcal{F}|, \omega$
都不大于$\max\{|\Omega|,|\mathcal{L}|\}$,
根据\href{../../01 一阶语言/01 一阶语言#nameddest=sec:定理A1.1}{引理}即得.
\hfill $\qedsymbol$

\vspace{10pt}

\hypertarget{sec:定理1.3}{}
\noindent\textbf{定理1.3}

给定一阶语言$\mathcal{L}$, $\mathcal{L}$-结构$B$.
对$B$上任意的变元赋值$\sigma$和非空的$W\subset\mathcal{V}$, 有
\[
    \langle\,\sigma[W]\,\rangle=\{t^{(B,\sigma)}\mid
    t\in T^{\mathcal{L}}\land\mathrm{var}(t)\subset W\}.
\]
特别地, 取$W=\mathcal{V}$就有
\[
    \langle\,\mathrm{Ran}\,\sigma\,\rangle=
    \{t^{(B,\sigma)}\mid t\in T^{\mathcal{L}}\}.
\]

\noindent{\kaishu 证明.}

第一, 证明
$\langle\,\sigma[W]\,\rangle\subset\{t^{(B,\sigma)}\mid
t\in T^{\mathcal{L}}\land\mathrm{var}(t)\subset W\}$.

首先, 对任意的$b\in\sigma[W]$, 存在$x\in W$使得$x^\sigma=b$,
并且$\mathrm{var}(x)=\{x\}\subset W$;

对任意的$b\in\{c^B\mid c\in\mathcal{C}\}$, 存在$c\in\mathcal{C}$使得
$c^B=b$, 并且$\mathrm{var}(c)=\varnothing\subset W$.

其次, 对任意的$n$元函数符号$f$, 假设$b_1,\cdots,b_n\in\{t^{(B,\sigma)}\mid
t\in T^{\mathcal{L}}\land\mathrm{var}(t)\subset W\}$, 即
\[
    \forall i,\,\exists t_i\in T^{\mathcal{L}}:\quad
    \mathrm{var}(t_i)\subset W\quad\land\quad b_i=t_i^{(B,\sigma)}.
\]
令$t=ft_1\cdots t_n$, 则$\mathrm{var}(t)\subset W$并且
$f^B(b_1,\cdots,b_n)=t^{(B,\sigma)}$.

第二, 对所有项$t$归纳证明
$\mathrm{var}(t)\subset W\to t^{(B,\sigma)}\in\langle\,\sigma[W]\,\rangle$.

首先, 如果$t$是变元且$\mathrm{var}(t)\subset W$, 则$t\in W$, 从而
$t^{(B,\sigma)}\in\sigma[W]\subset\langle\,\sigma[W]\,\rangle$;

如果$t$是常元, 则一定有$t^{(B,\sigma)}=t^B\in\langle\,\sigma[W]\,\rangle$.

其次, 对任意的$n$元函数符号$f$, 假设$t_1,\cdots,t_n$满足命题.
再假设$t=ft_1\cdots t_n$满足$\mathrm{var}(t)\subset W$.
此时对每个$t_i$也有$\mathrm{var}(t_i)\subset W$,
于是$t_i^{(B,\sigma)}\in\langle\,\sigma[W]\,\rangle$. 从而
$t^{(B,\sigma)}\in\langle\,\sigma[W]\,\rangle$. \hfill $\qedsymbol$





\newpage

\hypertarget{sec:定理1.4}{}
\noindent\textbf{定理1.4}

给定一阶语言$\mathcal{L}$, $\mathcal{L}$-结构$A,B$和它们上的变元赋值
$\sigma,\tau$, 以及非空的$W\subset\mathcal{V}$.

如果对任意的原子公式$\varphi$有
\[
    \mathrm{FV}(\varphi)\subset W\to
    ((A,\sigma)\models\varphi\leftrightarrow(B,\tau)\models\varphi),
\]
那么存在同构$p:\langle\,\sigma[W]\,\rangle\cong\langle\,\tau[W]\,\rangle$,
同时$p\circ\sigma\!\restriction_W=\tau\!\restriction_W$.

\noindent{\kaishu 证明.}

根据定理1.3即得
\[
    \langle\,\sigma[W]\,\rangle=\{t^{(A,\sigma)}\mid
    t\in T^{\mathcal{L}}\land\mathrm{var}(t)\subset W\},\quad
    \langle\,\tau[W]\,\rangle=\{t^{(B,\tau)}\mid
    t\in T^{\mathcal{L}}\land\mathrm{var}(t)\subset W\}.
\]
此时定义
\vspace{-3pt}
\[
    p=\{(t^{(A,\sigma)},t^{(B,\tau)})\mid
    t\in T^{\mathcal{L}}\land\mathrm{var}(t)\subset W\},\\[-3pt]
\]
易证$p$的定义域和值域就是$\langle\,\sigma[W]\,\rangle$和$\langle\,\tau[W]\,\rangle$.
接下来证明$p$是双射.

假设有$(a,b_1),(a,b_2)\in p$. 依定义, 存在变元集包含于$W$的项$t_1,t_2$使得
\[
    a=t_1^{(A,\sigma)}=t_2^{(A,\sigma)},\quad
    b_1=t_1^{(B,\tau)},\quad b_2=t_2^{(B,\tau)}.
\]
注意到$(A,\sigma)\models t_1\equiv t_2$,
于是$(B,\tau)\models t_1\equiv t_2$, 即$b_1=b_2$.

再假设有$(a_1,b),(a_2,b)\in p$, 同理易证$a_1=a_2$.
综上即得$p:\langle\,\sigma[W]\,\rangle\to\langle\,\tau[W]\,\rangle$是双射.

对任意的变元$x\in W$, 依定义即得$p(x^\sigma)=x^{\tau}$. 最后证明$p$是同态.

\begin{enumerate}
    \item 对任意的常元$c$, 依定义有$p(c^A)=c^B$.
    
    \item 对任意的$n$元函数符号$f$和$a_1,\cdots,a_n\in\langle\,\sigma[W]\,\rangle$,
    对每个$a_i$取项$t_i$使得
    \[
        \mathrm{var}(t_i)\subset W\quad\land\quad t_i^{(A,\sigma)}=a_i,
    \]
    则依定义有$p(a_i)=t_i^{(B,\tau)}$.

    \hspace{2.17em}
    令$t=ft_1\cdots t_n$, 则$\mathrm{var}(t)\subset W$, 于是
    \[
        p(f^A(a_1,\cdots,a_n))=p(t^{(A,\sigma)})
        =t^{(B,\tau)}=f^B(p(a_1),\cdots,p(a_n)).
    \]

    \item 对任意的$n$元关系符号$r$和$a_1,\cdots,a_n\in\langle\,\sigma[W]\,\rangle$,
    对每个$a_i$取项$t_i$使得
    \[
        \mathrm{var}(t_i)\subset W\quad\land\quad t_i^{(A,\sigma)}=a_i,
    \]
    则依定义有$p(a_i)=t_i^{(B,\tau)}$. 从而
    \begin{equation*}
        (a_1,\cdots,a_n)\in r^A\iff(A,\sigma)\models rt_1\cdots t_n\iff
        (B,\tau)\models rt_1\cdots t_n\iff(p(a_1),\cdots,p(a_n))\in r^B.
    \tag*{\qedsymbol}\end{equation*}
\end{enumerate}

\end{document}